% Exemplo: calendário de um mês com a biblioteca calendar do TikZ, com fins de semana em cinza e
% as datas de provas e entregas destacadas. Troque o mês nas datas de início e fim.
% Compilar: latexmk -pdf calendario.tex           Testado com: TikZ 3.1, TeX Live 2024 ou mais recente
%%BEGIN:calendariopreambulo
\documentclass[12pt, a4paper]{article}
\usepackage[T1]{fontenc}
\usepackage[brazilian]{babel}
\usepackage{tikz}
\usetikzlibrary{calendar}
%%END:calendariopreambulo

\begin{document}
\pagestyle{empty}

%%BEGIN:calendariomes
\begin{center}
\textbf{Novembro de 2026}\par\smallskip
\begin{tikzpicture}
  % cabeçalho: uma coluna por dia da semana, com o mesmo espaçamento dos dias
  \foreach \d [count=\i from 0] in {seg, ter, qua, qui, sex, sáb, dom}
    \node[font=\footnotesize\bfseries, text height=1.6ex, text depth=0.3ex]
      at (\i*8mm, 6mm) {\d};
  \calendar[dates=2026-11-01 to 2026-11-30, week list, day xshift=8mm,
            day yshift=6mm, day code={
              \node[anchor=center, minimum size=5.5mm, inner sep=0pt,
                    name=\pgfcalendarsuggestedname] {\tikzdaytext};}]
    if (weekend) [black!40]
    if (equals=2026-11-12) [nodes={draw=red, thick, circle}]            % prova 1
    if (equals=2026-11-26) [nodes={fill=blue!20, rounded corners=1pt}]   % entrega
  ;
\end{tikzpicture}\par\smallskip
\footnotesize\textcolor{red}{$\bigcirc$} prova \qquad
\colorbox{blue!20}{\phantom{0}} entrega do trabalho
\end{center}
%%END:calendariomes

\end{document}
